% page 588 du fac-simile (image p-587 du scan)
% bloc 152.4 x 246.3 mm ; marge gauche 8.0 mm
\pgc{-12.700mm}{0.000mm}{
\l{\cel{2}580}
\l{}
\l{\sou{texto}. I. La ipsa vorti di autoro, di lego, di akto (opo-}
\l{\cel{3}ze a la komenti). - II. Peci ek la libri santa quin}
\l{\cel{3}la predikisti citas, komence di sermono, di homilio,}
\l{\cel{3}kom aplikebla a la temo quan li intencas quik traktar.}
\l{\cel{3}DEFIRS.}
\l{}
\l{tezo. Propoziciono, asertata da ulu, e defensata da lu}
\l{\cel{3}kontre la personi qui kontestas olua exakteso, DEFIRS.}
\l{}
\l{tiaro.\cel{2}I. Kapovesto alta, uzata olim che la Oriento-populi.}
\l{\cel{3}- II. Kapovesto di la sacerdoto suprega, che la Hebrei.}
\l{\cel{3}- III. (\sou{che la katoliki)}. Kapovesto quan +weras la papo}
\l{\cel{3}dum la ceremonii, boneto kun, kom ornivi, tri kroni, e,}
\l{\cel{3}suprege, globo kun un krono. - DEFIRS.}
\l{}
\l{\sou{tibio}. (\sou{anat.}) La plu grosa ek la du osti di la gambo, avan}
\l{\cel{3}la peroneo. - EFIS.}
\l{}
\l{\sou{tifo}. (\sou{patol}.)Morbo kontagiema, quan karakterizas febro}
\l{\cel{3}sen-cesa, la stando malada di la mukozi, e c. - DEFIRS.}
\l{}
\l{\sou{tifoido}. (\sou{patol.}) Morbo infektiva quan produktas mikrobo}
\l{\cel{3}specifika (deskovrita da Eberth). - DEFIS.}
\l{}
\l{\sou{tifono}. (\sou{meteorol}.) Tempesto-vento, qua forportas kun su}
\l{\cel{3}- kun violento nerezistebla - to quon lu renkontras,}
\l{\cel{3}sur-tere e sur-maro. - EF.}
\l{}
\l{\sou{tigmotropismo}. (\sou{biol}.)Influo, da la kontakto, di la dire-}
\l{\cel{3}ciono che la tisui kreskanta. -}
\l{}
\l{\sou{tigro}. (\sou{zool.}) Animalo feroca, karnavida, ek la familio}
\l{\cel{3}\textquotedbl{}felini\textquotedbl{}, di qua la pilaro esas strioza e makuletoza.}
\l{\cel{3}L. felis tigris. - DEFIRS.}
\l{}
\l{\sou{tikar}. (\sou{netrans.}) I. (\sou{patol.}) (\sou{aludante la muskuli, preci}-}
\l{\cel{3}\sou{pue olti di la vizajo}).Kontraktesar konvulsatre, preske}
\l{\cel{3}kustumale. - II. (\sou{metaf.}) Agar gesto kustumale. - F.}
\l{}
\l{\sou{tiktakar}. (\sou{netrans.})\cel{2}(\sou{aludante, precipue, horlojo}).Pro-}
\l{\cel{3}duktar bruiseto bruska, ek du tempi intersucedanta qui}
\l{\cel{3}iteresas uniforme. - DEFS.}
\l{}
\l{\sou{til}. Prepoziciono qua indikas la arivo (en tempo, en spaco)}
\l{\cel{3}an termino quan onu ne transiras. (Kande la limito esas}
\l{\cel{3}ne punto, ma intervalo, lu sequesas da \textquotedbl{}exkluzite\textquotedbl{} od}
\l{\cel{3}\textquotedbl{}inkluzite\textquotedbl{} (segun la nociono). - E.}
\l{}
\l{\textquotedbl{}\sou{tilbury\textquotedbl{}}. Sorto di kabrioleto lejera, sen tendo. - EF.}
\l{}
\l{\sou{tildo}. (\sou{gram.)}\cel{2}Signo diakritika, qua havas ica formo : \textasciitilde{}}
\l{\cel{3}quan la Hispani lokizas sur la \textquotedbl{}n\textquotedbl{} kande li volas pro-}
\l{\cel{3}nuncigar lu \textquotedbl{}nyf\textquotedbl{} - e quan la Portugalani lokizas sur}
\l{\cel{3}\textquotedbl{}a\textquotedbl{} e \textquotedbl{}o\textquotedbl{} por pronuncigar li quale la F. \textquotedbl{}an\textquotedbl{} e \textquotedbl{}on\textquotedbl{},}
\l{\cel{3}(t.e. la E. \textquotedbl{}ang, ong\textquotedbl{}) rispektive. -}
}
